\documentclass[12pt]{article}  
\newcommand{\say}[1]{``#1''}  
\newcommand{\nsay}[1]{`#1'}  
\usepackage{endnotes}  
\newcommand{\B}{\backslash{}}  
\renewcommand{\,}{\textsuperscript{,}}  
\usepackage{setspace}   
\usepackage{tipa}  
\usepackage{hyperref}  
\begin{document}  
\doublespacing  
\section{\href{remembering-reconnect.html}{On Remembering to Reconnect}}  
First Published: 2026 June 21

\section{Draft 1: 21 June 2026}

As often happens when I start writing anything, I begin to remember all of the other forms of writing I do.  
I'd planned on restarting one web novel, beginning another, and starting to journal again.\footnote{though for now, journalling is a left handed activity}  
What I hadn't planned on, though, is the fact that some of the writing which feels most impactful to me is my communication with friends.

I have a variety of things about me which make it hard for me to reliably and recurrently\footnote{should be recurringly in my opinion} message those who matter to me.  
That's not a positive feature about me, but it's also not something that I think is likely to change.\footnote{and, I'm sure, it's linked to many of the positive things about me!}

Something that I have attempted in the past is making a list of people who I feel like I should message with some sort of regularity so that I remember to message them.  
This doesn't feel like a great solution, and not just because it didn't work.  
I also feel like it makes part of the connecting I do with some people more of a chore than something I do out of love.  
The distinction between doing something out of love and out of a sense of duty is not always clear, but it's also something that some people in my life have expressed dissatisfaction about in the past.

Why is this the focus of the post today?

It's the focus because today I remembered to reach out to a number of people I care deeply about today.  
I've not made a secret here or in life about the fact that these past months\footnote{years?} have been pretty rough for me mentally.  
As with many cycles in my behaviors, there is a clear feedback loop at play.  
When I feel bad, I am less likely to reach out to people.  
When I don't reach out to people, I both feel and become more isolated.  
When I am more isolated, I feel bad.  
This repeats until something breaks me out of the funk.

In contrast, I don't think that there's a virtuous\footnote{vicious and virtuous cycles are what I've a memory of being told behavior loops are} cycle in the same way for interacting with friends.  
Instead, I think that interacting with others serves mostly to put a floor on how bad I really feel.  
There's a large subjective difference between objective symptoms.  
As an example, if I'm in a lot of pain from a recent surgery, the pain doesn't necessarily recede from my attention or memory when someone laughs at my joke.  
What does happen, however, is that I am better able to ignore the sensation.  
Maybe that's a bad example.

I'm no more able to move my arm when people say hi.  
However, I am much less bothered by that fact, even when the limitations I have are made clearer.  
So, then, what are some ways that I might better keep in touch with those I love?\footnote{As I think I've said here and I know I've said in person a number of times, I think that one way we help with the male loneliness epidemic (as it's called) is by normalizing (if by force) the idea that men can and should love one another}.

I can update the lists and therefore have something in paper to reference to message people.

Pros of this approach:  
\begin{itemize}  
\item Makes it harder for one particular person to get lost in the shuffle  
\item Because it's something that happens on routine, neither of us has to feel like we're the one putting in a bunch of effort. Me, because I have a list that says to send a message, and them, because they're being reached out to  
\item I like lists  
\end{itemize}

Cons of this approach:  
\begin{itemize}  
\item Feels sterile. I'll never forget the holiday I sent greetings to a number of people, some of whom noted that I sent an identical message to others. By virtue\footnote{or, I guess, by vice, even though no one I have ever known or read uses such an expression} of the fact that I'm likely messaging many people at a time, the opening messages will become more rote  
\item Boom bust cycle to my messages. While I love being in contact with a number of people, I don't always love being in a large number of simultaneous conversations. When there's nothing else that I'm doing, it can be kind of nice.  
When I'm trying to also live my life, however, it feels less good to feel bombarded by a number of\footnote{admittedly self-inflicted} messages from a number of people.  
\item May be undesired from those around me. I'm sure that most of the people I don't message as much as I would love to also feel that it would be nice if we talked more. However, I'm almost positive that there are people who intentionally stopped responding to my messages or otherwise failed to continue a conversation.  
\item Similar to the first two points, the conversations by necessity cannot be quite as deep. I'm splitting my focus and my emotional energy in multiple directions when I have conversations with multiple people at a given time.  
\item Hard to maintain. Historically speaking, I have not done this.  
\end{itemize}

What's another idea?  
I guess I can keep on with what I'm doing, but that is of course not really a solution.

Third idea: I could keep a bunch of names in a jar and then pull one out whenever I'm feeling disconnected.

Pros of this concept:

\begin{itemize}  
\item I don't have the many names at once issue  
\item I don't have the \say{Hmmm... why do I get messages from J every X days} thought in the back of friends' minds  
\item I have a fun game I get to play!  
\end{itemize}

Cons of this concept:

\begin{itemize}  
\item As with the first, means that some people will likely feel I'm messaging too often.  
\item As with the first, requires me to remember two distinct actions: open the jar and also write a message from the jar.  
\item Still feels somehow non-genuine  
\end{itemize}

Message people on special events related to them. Big pro of this is that there's a reason that I'd be messaging them. Big con is they're probably overwhelmed by the sheer number of messages going out.

Another idea: maybe I just remake the daily to do list on this blog. Have an item for \say{have I messaged a friend to re-open lines of communication this (day? week? month?)}.  
Assuming that I can keep up with the blog\footnote{big assumption, almost certainly not something I'll be able to do given past performance}, that's just a lil reminder each day that I've got something to do.  
I think that's maybe the best idea.

Also, while I do feel like part of what some people\footnote{if only me in the future} want from this site is updates about my life, I don't feel like they go well in with many of the topics I'm discussing.  
Right now the solution has been a first draft that's much more stream of consciousness, but I don't think that is necessarily the best or most stable long term option.  
So, sure, let's add a new section to the bottom of each post\footnote{below the Upcoming Posts section} where I can do some reflections.

Let's do that daily reflection now!

Ok so going forward, this is likely a reflection I'd want to do \textit{before} I write the post, but then we get into the whole \say{how do I make that work what with the whole \nsay{I sometimes am working on multiple posts at once} thing}?

So where was i?

...

Oh, yeah I'm just going to use the daily reflection section of the blog to remind me to keep in touch with people.  
For now, I'm hoping that works.  
And at this exact moment, I'm feeling bathed in the love and good positive intentions of those who love me.

\section{Upcoming Posts}  
\begin{itemize}

\item Wonder  
\item Curse of knowledge, esp re. Music Theory\footnote{stolen from a prior list}  
\item Overview of a course on astrochem  
\item RebelFIt

\end{itemize}

\section{Daily Reflection: 21 June 2026}  
\begin{itemize}  
\item How's things?

Things are generally on an uphill swing. I think that I'm more and more realizing that my days are best when I start the day by doing something that nourishes me, like writing in a journal or moving.

\item What's a recent win?

I got the blog post done two days ago, have a blog post today, feel generally good, and practiced left-handed penmanship today!  
Biggest recent win is probably that I convinced my brother and father to read the last post I made.

\item You have the big tapestry embroidery project! How's that going?

Still have yet to put any real stitches on, but have done some more design work. At this point, I really think that I need to focus on the top and bottom layers, because I feel comfortable enough with the woven layer that I don't think that I'll really need too too much work with adapting it for any of the other questions.  
I do want to get better at stitching with my left hand. I keep defaulting to the right, which hurts a lot, and that's probably very not good for me.

\item Do you feel like you're taking good care of yourself?

Kind of! Yesterday wasn't the best, because I didn't really eat or drink as much as I should have, I didn't move, and I didn't do the activities which feed me.\footnote{other than, of course, the embroidery design, which is just very fun.}  
I'm generally feeling like today is going better, but this reflection is making me want to grab some water and a quick stretch before Father's Day Festivities.\footnote{... that would have been perhaps a more fitting post for today. Alas}

\item What could you be doing better?

Ooh, this question is definitely a great meta question. When I'm not doing well, the answer is of course going to feel like everything, which means that I can look at the answer and my immediate default to see if I need to start looking for something.  
I can be better about following through on the internal commitments I made.  
I told myself that while I'm recovering, there were some activities I would try doing to force myself to feel better, and I've not been the most consistent about them.

\item Are you appropriately keeping in touch with those you love?

As of today, mostly yeah.  
In general, not so much.  
I wish that I knew which people in my life liked my being a nuisance\footnote{there has to be a positive word for this idea, but like which of my friends go \say{ope! got another message, glad I'm still in thoughts even when I forget to write back}} and which would prefer me take a hint when they don't respond.   
Alas, this is real life, and I'm not brave enough to ask people these questions.

I've at the very least re-opened communication lines with people, so there's a chance that the communications are going to be better in coming days!

I also have the weekly brother call this evening, which will be really nice.

My schedule didn't line up with my partner's yesterday, so we didn't have the chance to call, but I'm optimistic that they'll line up today.\footnote{read: I'm fairly sure I have no plans today that are going to be urgent and important enough that I can't push them aside for the length of a call}

I miss writing letters, but that's probably an activity that's generally best saved for when I have my mobility.

\item Big upcoming events? Preparation for them?

\begin{itemize}

\item I have a two week camp equivalent coming up in mid July. I need to fill out some paperwork for that and need to figure out exactly where to send it.  
\item I'm visiting my partner in mid July. I need to book tickets for that

\end{itemize}

\item Once again, then, how's it going?

I think that it's going pretty well.  
The slope is absolutely trending upwards and fairly steeply, which is really nice to have be the case.  
I'm going to take a little bit of time to stretch and drink some water\footnote{even though I would normally also fill that time with random background noise, I'm going to use the time to instead sit and think quietly or at least not be deafening myself with noise. I'm realizing that quiet time, sucky as it often is, is usually very good for me. Is this really best as a footnote? Maybe or maybe not. Regardless, I think that most who read my blog are the sort who also read the footnotes, so really it's probably half of one and six dozen of another (wait. no, the opposite. Half dozen or six).}.  
Or maybe not. It may also be Father's Day festivities.\footnote{nope! I inherited some of my toxic \say{every day is a great day for labor} from him}

\end{itemize}


\end{document}